\section{矢量分析习题解答}

\begin{enumerate}[label=\arabic*.]
% 第1题
\item 设有标量场 $\varphi = 2xy^2 - z^3$，求 $\varphi$ 在点 $M(2,-1,1)$ 处沿该点至 $P(3,1,-1)$ 方向的方向导数。在点 $(2,-1,1)$ 沿什么方向导数达到最大值？其值是多少？

\textbf{解答：}
\begin{enumerate}[label=(\arabic*)]
\item 先求梯度 $\nabla\varphi$：
\[
\nabla\varphi = \frac{\partial\varphi}{\partial x}\vec{e}_x + \frac{\partial\varphi}{\partial y}\vec{e}_y + \frac{\partial\varphi}{\partial z}\vec{e}_z
= 2y^2\vec{e}_x + 4xy\vec{e}_y - 3z^2\vec{e}_z
\]
代入点 $M(2,-1,1)$：
\[
\nabla\varphi\big|_M = 2\vec{e}_x - 8\vec{e}_y - 3\vec{e}_z
\]

\item 方向矢量 $\vec{l} = \overrightarrow{MP} = (1,2,-2)$，单位化：
\[
\vec{e}_l = \frac{1}{3}\vec{e}_x + \frac{2}{3}\vec{e}_y - \frac{2}{3}\vec{e}_z
\]

\item 方向导数：
\[
\frac{\partial\varphi}{\partial l}\bigg|_M = -\frac{8}{3}
\]

\item 方向导数最大值出现在梯度方向，值为：
\[
|\nabla\varphi\big|_M| = \sqrt{77}
\]
\end{enumerate}

\textbf{结论：}
方向导数为 $\boxed{-\dfrac{8}{3}}$，最大值为 $\boxed{\sqrt{77}}$。

\vspace{1em}
% 第2题
\item 求标量场 $\varphi = x^3y^4z^2$ 的梯度场的散度。

\textbf{解答：}
\[
\nabla\varphi = 3x^2y^4z^2\vec{e}_x + 4x^3y^3z^2\vec{e}_y + 2x^3y^4z\vec{e}_z
\]
\[
\nabla\cdot(\nabla\varphi) = 6xy^4z^2 + 12x^3y^2z^2 + 2x^3y^4
\]

\textbf{结论：}
$\boxed{6xy^4z^2 + 12x^3y^2z^2 + 2x^3y^4}$。

\vspace{1em}
% 第3题
\item 若矢量 $\vec{A} = x^2\vec{e}_x + y^3\vec{e}_y + (3z - x)\vec{e}_z$，求(1) 散度；(2) 旋度。

\textbf{解答：}
\begin{enumerate}[label=(\arabic*)]
\item 散度：
\[
\nabla\cdot\vec{A}\big|_{(1,0,-1)} = 5
\]
\item 旋度：
\[
\nabla\times\vec{A} = \vec{e}_y
\]
\end{enumerate}

\textbf{结论：}
散度 $\boxed{5}$，旋度 $\boxed{\vec{e}_y}$。

\vspace{1em}
% 第4题
\item 已知电场强度 $\vec{E} = E_0\cos\theta\vec{e}_r - E_0\sin\theta\vec{e}_\theta$，求散度与旋度。

\textbf{解答：}
\[
\nabla\cdot\vec{E} = 0,\quad \nabla\times\vec{E} = \vec{0}
\]

\textbf{结论：}
$\boxed{\nabla\cdot\vec{E} = 0}$，$\boxed{\nabla\times\vec{E} = \vec{0}}$。

\vspace{1em}
% 第5题
\item 求矢量场对球面的通量。

\textbf{解答：}
由高斯公式：
\[
\nabla\cdot\vec{A} = 1
\]
\[
\iiint_V 1\,dV = \frac{4}{3}\pi a^3
\]

\textbf{结论：}
通量为 $\boxed{\dfrac{4}{3}\pi a^3}$。

\vspace{1em}
% 第6题
\item 求旋度与环量面密度。

\textbf{解答：}
\[
\nabla\times\vec{A}\big|_M = -\vec{e}_x - 3\vec{e}_y + 4\vec{e}_z
\]
\[
\mu_n = \frac{1}{3}
\]

\textbf{结论：}
旋度 $\boxed{-\vec{e}_x - 3\vec{e}_y + 4\vec{e}_z}$，环量面密度 $\boxed{\dfrac{1}{3}}$。

\end{enumerate}